%
% pendel.tex -- links das mathematische Pendel, rechts die Abweichung
% der exakten Schwingungsdauer von der Kleinwinkelnaeherung
% (Daten aus pendel-werte.dat).
%
\documentclass[tikz]{standalone}
\usepackage{amsmath}
\usepackage{times}
\usepackage{txfonts}
\usepackage{pgfplots}
\pgfplotsset{compat=1.18}
\usetikzlibrary{arrows,math}
%
% farben.tex -- Farben aus buch/common/macros.tex
%
\definecolor{darkgreen}{rgb}{0,0.6,0}
\definecolor{darkred}{rgb}{0.8,0,0}
\definecolor{orange}{rgb}{1,0.6,0}
\definecolor{gelb}{rgb}{1,0.8,0}

\begin{document}
\begin{tikzpicture}[>=latex, thick]

% ---------- linkes Panel: Pendel ----------
\def\Ls{2.60}
\def\th{35}

\fill (0,0) circle (0.06);
\node[anchor=south east] at (0.02,0.06) {$O$};

\draw[dashed, black!45] (0,0) -- (0,-\Ls-0.35);
\draw (0,0) -- ({\Ls*sin(\th)},{-\Ls*cos(\th)})
	node[midway, right=3pt] {$l$};

\draw[black!55] (0,-0.95) arc[start angle=-90, end angle=-90+\th, radius=0.95];
\node[black!55] at ({1.22*sin(\th/2)},{-1.22*cos(\th/2)}) {$\vartheta_0$};

\fill[darkred] ({\Ls*sin(\th)},{-\Ls*cos(\th)}) circle (0.13);
\node[darkred, anchor=west] at ({\Ls*sin(\th)+0.18},{-\Ls*cos(\th)}) {$m$};

\draw[->, black!55] ({\Ls*sin(\th)},{-\Ls*cos(\th)-0.28})
	-- ++(0,-0.62) node[anchor=north] {$g$};

% ---------- rechtes Panel: Abweichung der Schwingungsdauer ----------
\begin{axis}[
	at={(4.6cm,-3.55cm)}, anchor=south west,
	width=9.4cm, height=6.0cm,
	xlabel={$\vartheta_0$ in Grad},
	ylabel={Abweichung von $T_0$ in \%},
	xmin=0, xmax=170,
	ymin=0, ymax=150,
	xtick={0,30,60,90,120,150},
	ytick={0,25,50,75,100,125,150},
	grid=both,
	grid style={black!15},
]

\addplot[dashed, black!50, forget plot]
	coordinates {(0,18.03) (90,18.03) (90,0)};
\node[black!50, anchor=south west, font=\small] at (axis cs:4,20) {$18\,\%$};

\addplot[darkred, very thick, smooth] table[x=theta, y=abw] {pendel-werte.dat};

\end{axis}
\end{tikzpicture}
\end{document}
